\section{Global identification of the positive quadratic parameters}
\label{sec:v66-global-contact}

The preceding finite-jet calculation does not by itself exclude a second
positive quadratic branch.  We resolve this question directly, without
an analytic inverse theorem.  A Schur-complement identity converts the
curvature inverse into a strict contraction with constants determined
by the observed action Hessians.  This proves global identification of
the positive quadratic parameters, characterizes their image and gives
an a posteriori residual estimate.  Together with the signed higher-order
blocks, it supplies the finite-jet input for the actual smooth inverse
in Section~\ref{sec:v68-smooth-contact}; quadratic identification alone
is not asserted to identify a complete smooth germ.

Throughout this section the polygon, graph coordinates and signs are
those of Section~\ref{sec:v65-periodic-inverse}.  In particular
$k_i=L_i^{-1}>0$, $c_i=F_i''(0)>0$ and $s_i=(S_i^-)''(0)$;
indices are modulo the physical contact period $r$.  We regard $k$ as
fixed and write $\mathcal R_k(c)=s$.  This is a map of quadratic data,
independent of the higher contact jets.  Its definition uses the actual
half-line construction, not an arbitrary solution of a Riccati equation.

\subsection{The quadratic inverse}

\begin{lemma}[Boundary Schur complement]\label{lem:v66-riccati}
For every positive curvature vector $c$, the actual half-line action
Hessians satisfy $0<c_i<s_i$ and
\begin{equation}\label{eq:v66-schur}
 s_i=k_i+c_i-\frac{k_i^2}{k_i+c_{i+1}+s_{i+1}},\qquad
 a_i=\frac{k_i}{k_i+c_{i+1}+s_{i+1}}\in(0,1).
\end{equation}
Here $a_i$ is the unsigned one-step contraction in
Lemma~\ref{lem:v65-signed-jacobi}.  Conversely, positive vectors $c,s$
satisfying the first identity have $s=\mathcal R_k(c)$.
\end{lemma}
\begin{proof}
The half-line starting at $i$ is obtained by adjoining its first edge
to the half-line starting at $i+1$.  At quadratic order its stationary
value is the minimum over $y$ of
\[
 \tfrac12(k_i+c_i)u^2-\epsilon_i k_iuy
                  +\tfrac12(k_i+c_{i+1}+s_{i+1})y^2.
\]
The strictly positive half-line Hessian established in the preceding
section gives $s_{i+1}>0$.  Eliminating $y$ proves the first identity;
the stationary ratio is $y/u=\epsilon_i a_i$.  The denominator exceeds
$k_i$, proving both $a_i<1$ and $s_i>c_i$.

For the converse, define $a_i$ by the second identity.  Substituting the
first identity at $i+1$ gives
\[
 \frac{k_i}{a_i}
   =k_i+k_{i+1}+2c_{i+1}-k_{i+1}a_{i+1}.
\]
Thus the products $y_j=a_i\cdots a_{i+j-1}u$ solve the sign-conjugated
Jacobi equation with boundary value $u$.  They decay exponentially
because there are finitely many phases and every $a_i<1$.  Coercivity
makes this decaying solution unique: the difference of two such solutions
has zero boundary value, and summation of its positive Jacobi quadratic
form gives zero.  Restoring the signs gives the actual linearized
half-line.  Its initial action Hessian is $k_i+c_i-k_i a_i=s_i$.
Consequently these are the actual action Hessians.
\end{proof}

For $s\in(0,\infty)^r$ and $z\in[0,\infty)^r$, put
\begin{equation}\label{eq:v66-projection-map}
 \begin{aligned}
 (\Phi_s z)_i&=
 \left[s_i-k_i+\frac{k_i^2}{k_i+s_{i+1}+z_{i+1}}\right]_+,
 \qquad [t]_+=\max(t,0),\\
 q(s)&=\max_i\left(\frac{k_i}{k_i+s_{i+1}}\right)^2<1.
 \end{aligned}
\end{equation}
The positive-part operation specifies an inverse algorithm even for
positive data that are not realizable by strictly positive curvatures.
It does not add zero-curvature tables to the geometric class.

\begin{theorem}[Global curvature coordinates]\label{thm:v66-curvature}
The map $\mathcal R_k:(0,\infty)^r\longrightarrow(0,\infty)^r$
is one-to-one and a real analytic diffeomorphism onto its open image.
That image and the inverse admit the following description.

For every $s>0$, the map $\Phi_s$ has a unique fixed point $z(s)$ in
the closed positive orthant.  Its iterates from any $z^{(0)}\ge0$
converge to $z(s)$, with
\begin{equation}\label{eq:v66-iteration-error}
 \|z^{(m)}-z(s)\|_\infty
 \le \frac{q(s)^m}{1-q(s)}
                    \|z^{(1)}-z^{(0)}\|_\infty\quad(m\ge0).
\end{equation}
The first-increment estimate \eqref{eq:v66-iteration-error} is an
\emph{a priori} bound.  A residual at the current iterate gives the
\emph{a posteriori} bound
\begin{equation}\label{eq:v67-residual-error}
 \|z^{(m)}-z(s)\|_\infty\le E_m(s),\qquad
 E_m(s):=\frac{\|\Phi_s z^{(m)}-z^{(m)}\|_\infty}{1-q(s)}.
\end{equation}
For exact iteration its numerator is
$\|z^{(m+1)}-z^{(m)}\|_\infty$.
Precisely,
\begin{equation}\label{eq:v66-image}
 s\in\mathcal R_k((0,\infty)^r)
       \quad\Longleftrightarrow\quad z_i(s)>0\text{ for every }i;
 \qquad \mathcal R_k^{-1}(s)=z(s).
\end{equation}
For any $s,\widetilde s>0$, define the observable constant
\[
 q_* = \max_i
 \left(\frac{k_i}{k_i+\min(s_{i+1},\widetilde s_{i+1})}\right)^2<1.
\]
Then, including for data outside the geometric image,
\begin{equation}\label{eq:v66-pair-lipschitz}
 \|z(s)-z(\widetilde s)\|_\infty
       \le\frac{1+q_*}{1-q_*}\|s-\widetilde s\|_\infty.
\end{equation}
\end{theorem}
\begin{proof}
For $z\ge0$ the expression inside the positive part is strictly smaller
than $s_i$.  Hence $\Phi_s$ maps the closed orthant into
$\prod_i[0,s_i]$.  Its untruncated expression has derivative in
$z_{i+1}$ equal to
$-k_i^2/(k_i+s_{i+1}+z_{i+1})^2$, of absolute value at most $q(s)$.
The positive part is one-Lipschitz.  The map is therefore a strict
contraction in the maximum norm on the complete closed orthant.
The contraction theorem and summation of the successive differences
prove existence, uniqueness and \eqref{eq:v66-iteration-error}.
For any $v\ge0$, the triangle inequality and contraction give
\[
 \|v-z(s)\|_\infty
 \le\|v-\Phi_s v\|_\infty+q(s)\|v-z(s)\|_\infty.
\]
This proves \eqref{eq:v67-residual-error}.  It also provides a bound
for an inexact evaluation: if $w$ satisfies
$\|w-\Phi_s v\|_\infty\le\eta$, replace the residual numerator
by $\|w-v\|_\infty+\eta$.  An error bound for evaluating the map,
when used, is part of this certificate and is not assumed to be zero.

If $s=\mathcal R_k(c)$ for $c>0$, Lemma~\ref{lem:v66-riccati} says
that $c=\Phi_s c$.  Therefore $c=z(s)$, which proves global
injectivity.  If $z(s)>0$, no positive part is active at the fixed
point, so \eqref{eq:v66-schur} holds with $c=z(s)$.  The converse
in that lemma proves the other implication in \eqref{eq:v66-image}.
This proves the image characterization without identifying formal
positive data with geometric data.

For the two-point estimate, changing $s_i$ contributes at most
$\|s-\widetilde s\|_\infty$ to $\Phi_s z$; changing $s_{i+1}$
in its denominator contributes at most
$q_*\|s-\widetilde s\|_\infty$.  All intermediate denominators
are bounded below by the one used in $q_*$.  Changing $z$ contributes
at most $q_*\|z-\widetilde z\|_\infty$.  Apply these bounds at
the two fixed points and move the last term to the left.  This proves
\eqref{eq:v66-pair-lipschitz}.

On the image the fixed point is strictly positive.  Locally the positive
part can be omitted, and the remaining equation is real analytic in
$(s,c)$.  Its derivative in $c$ is $I+T_2$, where
$(T_2 v)_i=a_i^2v_{i+1}$, and has a convergent Neumann inverse because
$\|T_2\|_\infty<1$.  Thus the inverse is locally real analytic.
The forward half-line Hessian is real analytic in $c$: for example the
positive half-line Jacobi operator with periodic coefficients has a locally
uniformly convergent
resolvent expansion, and the boundary Schur complement is its analytic
matrix element.  Differentiating \eqref{eq:v66-schur} also gives
\begin{equation}\label{eq:v66-differential}
 (I-T_2)\,ds=(I+T_2)\,dc.
\end{equation}
Both factors are invertible.  Hence the forward map is a local real
analytic diffeomorphism at every $c>0$.  Its image is open, and the
already proved injectivity makes these local inverses one global inverse
on that image.
\end{proof}

The projection in \eqref{eq:v66-projection-map} is useful only in the
reconstruction algorithm; it plays no role in the analytic inverse near
a realized datum.  Formula~\eqref{eq:v66-iteration-error} also gives a
finite sufficient test for membership in the image: if every component
of an iterate exceeds either certified error bound, then $z(s)>0$.
For the residual bound this test is $\min_i z_i^{(m)}>E_m(s)$.
It certifies membership in the quadratic image, not extension to a
prescribed collection of globally disjoint analytic obstacles.
The constants are allowed to deteriorate as an action Hessian tends to
zero.  No uniform assertion at a grazing or zero-curvature limit is used.

There is a second direct uniqueness check.  If two positive curvature
vectors have the same $s$, subtraction of \eqref{eq:v66-schur} gives
\[
 c_i-\widetilde c_i
       =-a_i\widetilde a_i(c_{i+1}-\widetilde c_{i+1}).
\]
Iteration around the cycle forces the difference to vanish, since the
product of the $a_i\widetilde a_i$ is strictly smaller than one.
This argument treats odd and even physical periods alike.  It uses only
quadratic data; the signs in odd higher contact orders still have to be
retained in the geometric reconstruction.

For a normal two-contact channel, $k_0=k_1=1/g$ and $c_i$ are the
physical curvatures.  Set $H=(1+g^2s_0s_1)^{1/2}$.  The recovered
curvatures are
\[
 c_0=\frac{H\sqrt{s_0/s_1}-1}{g},\qquad
 c_1=\frac{H\sqrt{s_1/s_0}-1}{g}.
\]
These are precisely the alternating curvature formulas of
Theorem~\ref{thm:v26-density-inverse}.  Indeed put
$u_i=gs_i$ and $C_i=1+gc_i$; then $C_0C_1=1+u_0u_1$ and
$C_0/C_1=u_0/u_1$, which give
$(C_0-u_0)(C_1+u_1)=1$ and its reversed version.
These are the two identities in \eqref{eq:v66-schur}.
Whenever the displayed curvatures are positive, uniqueness identifies
them with the fixed point.  This is a specialization of the geometric
inverse, not a reduction of the two observation models.

Not every positive vector is an action datum for positive curvatures.
For example $r=2$, $k_0=k_1=1$ and $s=(1/10,10)$ give
$z(s)=(0,109/11)$.  The first component of the untruncated map is
negative for every $z\ge0$, and the second fixed-point equation gives
the stated value.  Thus this $s$ is excluded by \eqref{eq:v66-image},
even though the stable reconstruction algorithm is defined there.

